\section{Stability Controls and Policy Validation}

Because the controller operates with partial and periodically refreshed
statistics, stability is enforced at two levels.  During learning, replay,
target-network synchronization, bounded rewards, and epsilon decay limit the
rate at which value estimates change.  During execution, feasibility masks,
minimum-improvement thresholds, cooldowns, and the receipt-gated handoff protocol
prevent a high Q-value from bypassing required operating rules. These controls
bound the available transitions. Their role is operational stability, distinct
from a proof of learned-policy convergence.

Classical tabular Q-learning guarantees assume a stationary, fully observable
Markov decision process and sufficient visitation of every state-action pair.
The client instead receives delayed partial observations in an environment
whose queue arrivals, acknowledgement capacity, and competing clients can
change. Policy quality is therefore evaluated under the stated workloads,
without a global-optimality guarantee. The finite v3 traces describe training
behavior rather than Bellman-loss convergence. The v4--v6 results distinguish
training completion from action coverage. The validation protocol uses
action-coverage, fault-safety, churn, and performance gates to assess each
frozen policy.

RL evaluation also separates adaptation from regression. Biagiola and Tonella
test how RL systems adapt to changed environments while retaining prior
capabilities \cite{Biagiola2022Plasticity}. Patterson et al. discuss performance
variation, baseline construction, hyperparameter selection, and experimenter
bias in empirical RL \cite{Patterson2024EmpiricalDesign}. These concerns motivate
reporting the training, frozen-policy comparison, and live handoff evidence
separately. A successful handoff and a higher comparative reward answer
different evaluation questions.

\section{Discussion}

The main systems finding is that subscriber capacity and client access to that
capacity are separate. Adding Pods supplies service capacity; moving existing
subscription-bearing connections makes it reachable by persistent clients.
The guarded-greedy and DQL policies both achieve this in the scale-out model.
Their nearby throughput results support the shared mechanism rather than
requiring a claim that learning beats a strong greedy rule in every workload.

Policy choice depends on the operational objective. Greedy control gives a
direct, inspectable response to measured pressure. DQL can account for delayed
costs and several competing objectives. A subscription move changes future
backlog distribution; migration trades immediate handoff cost for later
capacity. Both selectors operate above the same bounded, checked actions.

The action-contract studies provide a further design lesson: a useful
executor and a policy that reliably selects it are separate requirements.
The cap-restoring scale-up preflight improved reward and maximum subscription
count, but the expanded learner study did not pass every coverage gate.
Appendix~\suppref{app:action-diagnostics} retains these results. They explain why
the live checkpoint uses a compatible, restricted rebalancing profile.

\section{Conclusion}

We presented a probe-guided, client-side framework for rebalancing
persistent subscription connections without requiring cluster-wide client
visibility.  The mechanism combines per-subscription queues, active-connection
telemetry, and a load-balancer-assigned probe. Greedy rules or DQL select
bounded actions, and a shared executor carries out receipt-gated handoff.
The technical contribution is the link between partial observations,
endpoint-preserving probe promotion, and selective subscription reassociation.

The primary experiment demonstrates recovery after a one-to-four pod capacity
increase with 20,000 subscriptions and 775 million queued events. Guarded-greedy
migration and frozen DQL deliver 221,260 and 219,762 events/s, respectively,
compared with 58,000 events/s for unchanged placement. Frozen DQL
uses all four pods, delivers 3.79 times the mean throughput, and leaves 27.52\%
less backlog after 20 minutes than retaining the ten original connection
assignments. Its 12.30\% reduction in backlog area satisfies the registered
paired statistical criterion. The adverse-probe ablation complements this
result: evaluating the candidate before promotion reduces backlog area by
4.19\% relative to forcing one harmful migration.
The Oracle Cloud experiments establish live handoff execution, subscription
retention, and observed delivery behavior. The comparative throughput and
backlog-recovery gains are simulator measurements under the stated service
assumptions. These results support two evaluated policy options above the same
probe-guided handoff mechanism.
